\documentclass[11pt]{article}
\usepackage[margin=0.9in]{geometry}
\usepackage{amsmath,amssymb,amsthm}
\usepackage{booktabs,enumitem,xcolor,fancyhdr,array,longtable}
\usepackage[most]{tcolorbox}
\usepackage[hidelinks]{hyperref}

\definecolor{accent}{HTML}{1F4E79}
\definecolor{tipbg}{HTML}{FFF6DD}
\definecolor{exbg}{HTML}{EEF6EE}
\definecolor{methodclr}{HTML}{6A3D9A}

\newtcolorbox{keybox}[1][]{colback=accent!6, colframe=accent, boxrule=0.6pt,
  arc=2pt, left=6pt, right=6pt, top=4pt, bottom=4pt, fonttitle=\bfseries, #1}
\newtcolorbox{tipbox}[1][]{colback=tipbg, colframe=orange!70!black, boxrule=0.6pt,
  arc=2pt, left=6pt, right=6pt, top=4pt, bottom=4pt, fonttitle=\bfseries, #1}
\newtcolorbox{examplebox}[1][]{colback=exbg, colframe=green!40!black, boxrule=0.6pt,
  arc=2pt, left=6pt, right=6pt, top=4pt, bottom=4pt, fonttitle=\bfseries, breakable, #1}
\newtcolorbox{methodbox}[1][]{colback=methodclr!5, colframe=methodclr, boxrule=0.6pt,
  arc=2pt, left=6pt, right=6pt, top=4pt, bottom=4pt, fonttitle=\bfseries, title={Method}, #1}

\newcommand{\E}{\mathbb{E}}
\newcommand{\Var}{\operatorname{var}}
\newcommand{\Cov}{\operatorname{cov}}
\newcommand{\corr}{\operatorname{corr}}
\newcommand{\se}{\operatorname{se}}
\newcommand{\MSE}{\operatorname{MSE}}
\newcommand{\bias}{\operatorname{bias}}
\renewcommand{\P}{\mathbb{P}}
\newcommand{\given}{\mid}
\newcommand{\Bern}{\text{Bernoulli}}
\newcommand{\Norm}{\text{Normal}}
\newcommand{\SSR}{\mathrm{SSR}}
\newcommand{\SST}{\mathrm{SST}}
\newcommand{\Ybar}{\overline{Y}}
\newcommand{\Xbar}{\overline{X}}
\newcommand{\why}[1]{&&\text{\small\color{gray}(#1)}}
% move tags, e.g. \mv{3}
\newcommand{\mv}[1]{\tcbox[on line, colback=methodclr, colframe=methodclr, coltext=white, boxsep=0pt,
  left=2pt, right=2pt, top=1pt, bottom=1pt, arc=1.5pt]{\sffamily\scriptsize\bfseries M#1}}
\newcommand{\goal}[1]{\textbf{Goal:} #1\par}
\newcommand{\start}[1]{\textbf{Start from:} #1\par}
\newcommand{\moves}[1]{\textbf{Moves:} #1\par}

\setlength{\parindent}{0pt}
\emergencystretch=3em
\setlength{\parskip}{5pt}
\setlength{\headheight}{14pt}
\pagestyle{fancy}
\fancyhf{}
\lhead{Math 50 -- Midterm 2: deriving the formula sheet}
\rhead{Unit 2.4, Unit 3, Unit 4}
\cfoot{\thepage}

\title{\textbf{Math 50 Midterm 2: How to Derive Every Formula on the Sheet}\\[2pt]
\large A companion to \texttt{midterm2\_formulas.pdf} and \texttt{midterm2\_notes.pdf}}
\author{}
\date{}

\begin{document}
\maketitle
\thispagestyle{fancy}

\begin{tipbox}[title={How to use this document}]
Every formula on the formula sheet can be derived from the model equation with a small number of \emph{moves}. Section~\ref{sec:moves}
lists those moves (M1--M11) and translates each arrow on the formula sheet into a move. Every later derivation starts with a purple
\textbf{Method} box saying the goal, what to start from, and which moves to use. If you can reproduce the Method box from memory, you can
usually reproduce the algebra. The full steps follow, with the reason for each step in gray on the right, as in the notes.

Blue boxes hold the final result. Sections marked \emph{(extra)} are not on the sheet itself but come from the notes and use the same
methods. The last page is a cover-and-recall table: hide the right-hand columns and try to say the method for each formula.

Everything here uses only tools from the Unit 2--4 notes and the lectures (linearity of $\E$, bilinearity of $\Cov$, variance rules, the CLT,
and sums). Nothing new is introduced.
\end{tipbox}

\tableofcontents

%==========================================================
\clearpage
\section{The moves, and how to read the formula sheet's arrows}\label{sec:moves}
%==========================================================
\subsection{The eleven moves}
\renewcommand{\arraystretch}{1.35}
\begin{longtable}{@{}>{\raggedright\arraybackslash}p{0.9cm}>{\raggedright\arraybackslash}p{3.6cm}>{\raggedright\arraybackslash}p{11.3cm}@{}}
\toprule
& \textbf{Move} & \textbf{What you do / the rule it uses}\\
\midrule
\endhead
\mv{1} & Substitute the model & Replace $Y$ by $\beta_0+\beta_1X_1+\epsilon$ (or the two-predictor version). Almost every population derivation starts here.\\
\mv{2} & Take $\E[\cdot]$ & Linearity: $\E[aX+bY+c]=a\E X+b\E Y+c$. Use $\E[\epsilon]=0$.\\
\mv{3} & Take $\Cov(X_1,\cdot)$ & Bilinearity: $\Cov(X_1,aY+bZ+c)=a\Cov(X_1,Y)+b\Cov(X_1,Z)$. Constants drop out. $\Cov(X_1,X_1)=\Var X_1$. $\Cov(X_1,\epsilon)=0$ because the noise is independent of the predictors.\\
\mv{4} & Take $\Var(\cdot)$ & $\Var(aX+b)=a^2\Var X$. Independent pieces: variances add. In general $\Var(X+Y)=\Var X+\Var Y+2\Cov(X,Y)$.\\
\mv{5} & Condition & ``Given $X$'' means $X$ is a known constant. $\Var(Y\given X)$ is the variance of whatever is \emph{still random}.\\
\mv{6} & Rescale / standardize & Divide by standard deviations: $Z=(X-\mu)/\sigma$. Changes units, not shape.\\
\mv{7} & Average of $N$ iid pieces & $\E[\Ybar]=\mu$, $\Var(\Ybar)=\sigma^2/N$. A sum of independent Normals is Normal; otherwise use the CLT for large $N$.\\
\mv{8} & Plug in & Replace each expectation by a sample average, $\E[\cdot]\to\frac1N\sum_i$, and $\E X,\E Y$ by $\Xbar,\Ybar$.\\
\mv{9} & Normal $\to$ probability & Standardize, then use $\P(|Z|<1.96)=0.95$ and $\P(Z<z)=\Phi(z)$.\\
\mv{10} & Add and subtract & Insert $\pm\mu$ or $\pm\E[\hat\theta]$ inside a square, expand $(a+b)^2$, and show the cross term is zero.\\
\mv{11} & The sum trick & $\sum_i(X_i-\Xbar)=0$, so $\sum_i c\,(X_i-\Xbar)=0$ for any constant $c$ (including $c=\Ybar$).\\
\bottomrule
\end{longtable}

\subsection{Each arrow on the sheet is a move}
The labels on the arrows of the formula sheet are instructions. Here is what each one means and where it is done below.
\begin{longtable}{@{}>{\raggedright\arraybackslash}p{3.3cm}>{\raggedright\arraybackslash}p{4.2cm}>{\raggedright\arraybackslash}p{6.2cm}>{\raggedright\arraybackslash}p{1.6cm}@{}}
\toprule
\textbf{Arrow label} & \textbf{From $\to$ to} & \textbf{What you actually do} & \textbf{Section}\\
\midrule
\endhead
take $\Cov(X_1,\cdot)$ and $\E[\cdot]$ & one predictor $\to$ coefficients & \mv{1}\,\mv{3} then divide by $\Var X_1$; \mv{1}\,\mv{2} then solve for $\beta_0$ & \ref{sec:beta1}, \ref{sec:beta0}\\
take $\Var(\cdot)$ & one predictor $\to$ variance & \mv{1}\,\mv{4}, noise independent of $X_1$ & \ref{sec:varY}\\
rescale by $\sigma_{X_1}/\sigma_Y$ & coefficients $\to$ correlation & \mv{6}: multiply $\beta_1$ by $\sigma_{X_1}/\sigma_Y$ & \ref{sec:rho}\\
$\Var(Y\given X_1)=\sigma_\epsilon^2$ & variance $\to$ variance explained & \mv{5}: given $X_1$, only $\epsilon$ is random & \ref{sec:condvar}\\
square & correlation $\to$ variance explained & square $\rho$, use the variance equation & \ref{sec:rho2}\\
omit $X_2$ & two predictors $\to$ one predictor & the one-predictor slope of the \emph{two}-predictor world is $\beta_1'$ & \ref{sec:ovb}\\
$\Var(Y\given X_1,X_2)=\sigma_\epsilon^2$ & two predictors $\to$ variance explained & \mv{5} with both predictors known & \ref{sec:two-var}\\
$\beta_1\to\beta_1'$; $\ \beta_1,\beta_2$ & coefficients / two predictors $\to$ OVB & \mv{1}\,\mv{3} in the two-predictor model & \ref{sec:ovb}\\
estimate $\beta_1$ from $N$ observations & one predictor $\to$ sample dist.\ of $\hat\beta_1$ & \mv{8} to get $\hat\beta_1$, then \mv{11}\,\mv{2}\,\mv{4} & \ref{sec:ls}\\
$\hat\theta=\Ybar$ & sample distribution $\to$ sample mean & \mv{7} & \ref{sec:ybar}\\
$\hat\theta=\hat\beta_1$ & sample distribution $\to$ sample dist.\ of $\hat\beta_1$ & same recipe as $\Ybar$, with weights & \ref{sec:varb1}\\
(to bias/variance) & sample distribution $\to$ MSE & \mv{10} & \ref{sec:mse}\\
$\pm1.96\,\se$ & sample distribution $\to$ 95\% CI & \mv{9}, then rearrange the inequality & \ref{sec:ci}\\
null $\beta_1=0$ & sample dist.\ of $\hat\beta_1$ $\to$ hypothesis test & set $\beta_1=0$, \mv{6}\,\mv{9} & \ref{sec:test}\\
(CI $\to$ test) & 95\% CI $\to$ hypothesis test & chain of ``if and only if'' steps & \ref{sec:ci-test}\\
\bottomrule
\end{longtable}

\begin{keybox}[title={The general recipe for any sample distribution}]
Whenever you are asked ``what is the sample distribution of $\hat\theta$?'':
\begin{enumerate}[leftmargin=*,itemsep=1pt,label=\arabic*.]
\item Write $\hat\theta$ as a (weighted) sum of the data, e.g.\ $\Ybar=\sum_i\frac1NY_i$ or $\hat\beta_1=\sum_i w_iY_i$.
\item \textbf{Center:} take $\E[\cdot]$ \mv{2}. Compare with $\theta$ to get the bias.
\item \textbf{Spread:} take $\Var(\cdot)$ \mv{4}; independent data points, so variances add. The square root is $\se(\hat\theta)$.
\item \textbf{Shape:} a sum of independent Normals is Normal; otherwise the CLT says approximately Normal for large $N$ \mv{7}.
\end{enumerate}
\end{keybox}

\begin{tipbox}[title={If you blank on the exam}]
Write the model equation, then ask which arrow you need: for a \textbf{coefficient}, take $\Cov(X,\cdot)$; for an \textbf{intercept} or mean, take
$\E$; for a \textbf{spread}, take $\Var$ (squared coefficients, independent pieces add); for anything ``\textbf{given} $X$'', freeze $X$ as a constant;
for anything \textbf{estimated}, write it as a weighted sum of data and do center, spread, shape.
\end{tipbox}

%==========================================================
\clearpage
\section{One predictor: the population formulas}\label{sec:onepred}
%==========================================================
Throughout this section the model is
\[
Y=\beta_0+\beta_1X_1+\epsilon,\qquad \epsilon\sim\Norm(0,\sigma_\epsilon^2),\qquad \epsilon \text{ independent of } X_1 .
\]
Write $\sigma_{X_1}^2=\Var(X_1)$ and $\sigma_Y^2=\Var(Y)$. These are statements about the true (population) quantities, not data.

\subsection{The slope: $\beta_1=\Cov(X_1,Y)/\sigma_{X_1}^2$}\label{sec:beta1}
\begin{methodbox}
\goal{an expression for $\beta_1$ using only $X_1$ and $Y$.}
\start{the model equation.}
\moves{\mv{1} substitute the model inside $\Cov(X_1,Y)$; \mv{3} split it with bilinearity; solve for $\beta_1$.}
\end{methodbox}
\begin{align*}
\Cov(X_1,Y)&=\Cov(X_1,\ \beta_0+\beta_1X_1+\epsilon) \why{substitute the model}\\
&=\beta_1\Cov(X_1,X_1)+\Cov(X_1,\epsilon) \why{bilinearity; the constant $\beta_0$ drops out}\\
&=\beta_1\sigma_{X_1}^2+0 \why{$\Cov(X_1,X_1)=\Var X_1$; $\epsilon$ independent of $X_1$}
\end{align*}
Divide both sides by $\sigma_{X_1}^2$:
\begin{keybox}
\[\beta_1=\frac{\Cov(X_1,Y)}{\sigma_{X_1}^2}.\]
\end{keybox}
\textbf{Why this move works.} Taking $\Cov(X_1,\cdot)$ of the model ``kills'' everything not connected to $X_1$: constants vanish and the noise
is uncorrelated with $X_1$. Only the $\beta_1X_1$ term survives. You will use this same move for omitted-variable bias (Section~\ref{sec:ovb})
and the two-predictor system (Section~\ref{sec:system}).

\subsection{The intercept: $\beta_0=\E[Y]-\beta_1\E[X_1]$}\label{sec:beta0}
\begin{methodbox}
\goal{$\beta_0$.}
\start{the model equation.}
\moves{\mv{1}\,\mv{2} take the expectation of both sides; use $\E[\epsilon]=0$; solve for $\beta_0$.}
\end{methodbox}
\begin{align*}
\E[Y]&=\E[\beta_0+\beta_1X_1+\epsilon] \why{take $\E$ of both sides}\\
&=\beta_0+\beta_1\E[X_1]+\E[\epsilon] \why{linearity}\\
&=\beta_0+\beta_1\E[X_1] \why{$\E[\epsilon]=0$}
\end{align*}
\begin{keybox}
\[\beta_0=\E[Y]-\beta_1\E[X_1].\]
\end{keybox}
\textbf{Meaning.} The line passes through the point of means $(\E X_1,\E Y)$. The least squares version is $\hat\beta_0=\Ybar-\hat\beta_1\Xbar$.

\subsection{(extra) A binary predictor: $\beta_1$ is a difference of group means}\label{sec:binary}
\begin{methodbox}
\goal{read $\beta_1$ off group averages when $X_1\in\{0,1\}$.}
\start{the model, conditioned on each value of $X_1$.}
\moves{\mv{5} condition on $X_1=1$ and on $X_1=0$; \mv{2} take $\E$; subtract.}
\end{methodbox}
Given $X_1=x$, the terms $\beta_0+\beta_1x$ are constants, so $\E[Y\given X_1=x]=\beta_0+\beta_1x+\E[\epsilon]=\beta_0+\beta_1x$. Therefore
\[
\E[Y\given X_1=1]-\E[Y\given X_1=0]=(\beta_0+\beta_1)-\beta_0=\beta_1 .
\]
This is how the notes read slopes off plots (Lecture 7) and is the basis of the graders example and of the second OVB method in
Section~\ref{sec:ovb}.

\subsection{The variance equation: $\sigma_Y^2=\beta_1^2\sigma_{X_1}^2+\sigma_\epsilon^2$}\label{sec:varY}
\begin{methodbox}
\goal{split the variance of $Y$ into the part from $X_1$ and the part from noise.}
\start{the model equation.}
\moves{\mv{1}\,\mv{4} take $\Var$ of both sides; the two random pieces are independent, so their variances add.}
\end{methodbox}
\begin{align*}
\sigma_Y^2=\Var(Y)&=\Var(\beta_0+\beta_1X_1+\epsilon) \why{substitute the model}\\
&=\Var(\beta_1X_1)+\Var(\epsilon) \why{constant $\beta_0$ drops; $X_1,\epsilon$ independent}\\
&=\beta_1^2\sigma_{X_1}^2+\sigma_\epsilon^2 \why{$\Var(aX)=a^2\Var X$}
\end{align*}
\begin{keybox}
\[\sigma_Y^2=\beta_1^2\sigma_{X_1}^2+\sigma_\epsilon^2 .\]
\end{keybox}
\textbf{Common slip:} $\Var(\beta_1X_1)$ is $\beta_1^2\sigma_{X_1}^2$, not $\beta_1\sigma_{X_1}^2$. The coefficient comes out \emph{squared}.

\subsection{The conditional variance: $\Var(Y\given X_1)=\sigma_\epsilon^2$}\label{sec:condvar}
\begin{methodbox}
\goal{how much $Y$ varies once $X_1$ is known.}
\start{the model equation.}
\moves{\mv{5} once $X_1$ is known, $\beta_0+\beta_1X_1$ is a constant; \mv{4} a constant does not change variance.}
\end{methodbox}
\begin{align*}
\Var(Y\given X_1)&=\Var(\beta_0+\beta_1X_1+\epsilon\given X_1) \why{substitute the model}\\
&=\Var(\epsilon\given X_1) \why{$\beta_0+\beta_1X_1$ is a known constant}\\
&=\sigma_\epsilon^2 \why{$\epsilon$ independent of $X_1$}
\end{align*}
\textbf{Meaning.} $\sigma_\epsilon^2$ is the spread of the points around the line; $\sigma_Y^2$ is the total spread. Comparing the two is
exactly what $\rho^2$ does.

\subsection{Correlation: $\rho=\beta_1\,\sigma_{X_1}/\sigma_Y$}\label{sec:rho}
\begin{methodbox}
\goal{the slope in standard-deviation units.}
\start{either the definition $\rho=\Cov(X_1,Y)/(\sigma_{X_1}\sigma_Y)$ (method A) or the model in standardized variables (method B).}
\moves{A: substitute $\Cov(X_1,Y)=\beta_1\sigma_{X_1}^2$ from Section~\ref{sec:beta1} and cancel. \ B: \mv{1}\,\mv{6} standardize both
sides and read off the new slope.}
\end{methodbox}
\textbf{Method A (fast, use on the exam).}
\[
\rho=\frac{\Cov(X_1,Y)}{\sigma_{X_1}\sigma_Y}=\frac{\beta_1\sigma_{X_1}^2}{\sigma_{X_1}\sigma_Y}=\beta_1\frac{\sigma_{X_1}}{\sigma_Y}.
\]
\textbf{Method B (shows \emph{why} it is called ``rescaling'').} Let $Z_X=(X_1-\mu_X)/\sigma_{X_1}$ and $Z_Y=(Y-\mu_Y)/\sigma_Y$, with
$\mu_Y=\beta_0+\beta_1\mu_X$ from Section~\ref{sec:beta0}.
\begin{align*}
Z_Y=\frac{Y-\mu_Y}{\sigma_Y}
&=\frac{(\beta_0+\beta_1X_1+\epsilon)-(\beta_0+\beta_1\mu_X)}{\sigma_Y} \why{substitute $Y$ and $\mu_Y$}\\
&=\frac{\beta_1(X_1-\mu_X)+\epsilon}{\sigma_Y} \why{$\beta_0$ cancels}\\
&=\frac{\beta_1\sigma_{X_1}}{\sigma_Y}\cdot\frac{X_1-\mu_X}{\sigma_{X_1}}+\frac{\epsilon}{\sigma_Y} \why{multiply and divide by $\sigma_{X_1}$}\\
&=\underbrace{\frac{\beta_1\sigma_{X_1}}{\sigma_Y}}_{\text{new slope}}\,Z_X+\frac{\epsilon}{\sigma_Y}. \why{definition of $Z_X$}
\end{align*}
\begin{keybox}
\[\rho=\beta_1\frac{\sigma_{X_1}}{\sigma_Y}=\frac{\Cov(X_1,Y)}{\sigma_{X_1}\sigma_Y}
\qquad\text{= the slope of $Z_Y$ on $Z_X$.}\]
\end{keybox}

\subsection{Variance explained: $\rho^2=1-\Var(Y\given X)/\Var(Y)$}\label{sec:rho2}
\begin{methodbox}
\goal{$\rho^2$ as a fraction of variance.}
\start{the variance equation $\sigma_Y^2=\beta_1^2\sigma_{X_1}^2+\sigma_\epsilon^2$ (Section~\ref{sec:varY}).}
\moves{divide through by $\sigma_Y^2$; recognize $\beta_1^2\sigma_{X_1}^2/\sigma_Y^2=\rho^2$ (square Section~\ref{sec:rho}); replace
$\sigma_\epsilon^2$ by $\Var(Y\given X)$ (Section~\ref{sec:condvar}).}
\end{methodbox}
\begin{align*}
\sigma_Y^2&=\beta_1^2\sigma_{X_1}^2+\sigma_\epsilon^2 \why{variance equation}\\
1&=\frac{\beta_1^2\sigma_{X_1}^2}{\sigma_Y^2}+\frac{\sigma_\epsilon^2}{\sigma_Y^2} \why{divide by $\sigma_Y^2$}\\
1&=\rho^2+\frac{\sigma_\epsilon^2}{\sigma_Y^2} \why{$\rho^2=(\beta_1\sigma_{X_1}/\sigma_Y)^2$}\\
\rho^2&=1-\frac{\Var(Y\given X)}{\Var(Y)} \why{$\sigma_\epsilon^2=\Var(Y\given X)$, $\sigma_Y^2=\Var Y$}
\end{align*}
\begin{keybox}
\[\rho^2=1-\frac{\Var(Y\given X)}{\Var(Y)}=\frac{\beta_1^2\sigma_{X_1}^2}{\sigma_Y^2}.\]
\end{keybox}
\textbf{Bonus: why $-1\le\rho\le1$.} Both pieces $\rho^2$ and $\sigma_\epsilon^2/\sigma_Y^2$ are $\ge0$ and add to $1$, so $\rho^2\le1$.

%==========================================================
\clearpage
\section{Two (or more) predictors}\label{sec:twopred}
%==========================================================
Now the model is
\[
Y=\beta_0+\beta_1X_1+\beta_2X_2+\epsilon,\qquad\epsilon\sim\Norm(0,\sigma_\epsilon^2),\qquad\epsilon\text{ independent of }X_1,X_2 .
\]
$\beta_1$ is the effect of $X_1$ \emph{with $X_2$ held fixed}. We write $\beta_1'$ for the slope when $Y$ is regressed on $X_1$ alone.

\subsection{Conditional variance and $\rho^2$ with two predictors; adding a predictor never lowers $\rho^2$}\label{sec:two-var}
\begin{methodbox}
\goal{(a) $\Var(Y\given X_1,X_2)$; (b) show $\rho^2$ with both predictors $\ge$ $\rho^2$ with $X_1$ alone.}
\start{the two-predictor model.}
\moves{\mv{5} condition on what is known; \mv{4} the remaining random pieces are independent, so variances add; compare.}
\end{methodbox}
\textbf{(a)} Given both $X_1$ and $X_2$, the whole $\beta_0+\beta_1X_1+\beta_2X_2$ is a constant, so exactly as in Section~\ref{sec:condvar},
\[
\Var(Y\given X_1,X_2)=\Var(\epsilon)=\sigma_\epsilon^2,\qquad\text{so}\qquad \rho^2=1-\frac{\Var(Y\given X_1,X_2)}{\Var Y}=1-\frac{\sigma_\epsilon^2}{\sigma_Y^2}.
\]
(Taking $\Var$ of the model, with \mv{4}'s general rule since $X_1,X_2$ may be correlated:
$\sigma_Y^2=\beta_1^2\sigma_{X_1}^2+\beta_2^2\sigma_{X_2}^2+2\beta_1\beta_2\Cov(X_1,X_2)+\sigma_\epsilon^2$.)

\textbf{(b)} Given only $X_1$, the term $\beta_0+\beta_1X_1$ is a constant but $X_2$ and $\epsilon$ are still random:
\begin{align*}
\Var(Y\given X_1)&=\Var(\beta_2X_2+\epsilon\given X_1) \why{drop the known constant}\\
&=\beta_2^2\Var(X_2\given X_1)+\sigma_\epsilon^2 \why{$\epsilon$ independent of $X_1,X_2$: variances add}\\
&\ge\sigma_\epsilon^2=\Var(Y\given X_1,X_2). \why{a variance is $\ge0$}
\end{align*}
Since $\rho^2=1-\Var(Y\given\cdots)/\Var Y$, a smaller conditional variance means a bigger $\rho^2$:
\begin{keybox}
\[\rho^2_{\text{with }X_1,X_2}\ \ge\ \rho^2_{\text{with }X_1\text{ only}}.\]
Adding a predictor never decreases $\rho^2$ (and the same argument with sample quantities shows $R^2$ never decreases).
\end{keybox}

\subsection{Omitted-variable bias: $\beta_1'=\beta_1+\beta_2\beta_{1,2}$}\label{sec:ovb}
\begin{methodbox}
\goal{the slope $\beta_1'$ you get if you regress $Y$ on $X_1$ alone, when the truth has two predictors.}
\start{the one-predictor formula $\beta_1'=\Cov(X_1,Y)/\Var(X_1)$ (Section~\ref{sec:beta1}), which holds for \emph{any} $Y$.}
\moves{\mv{1} substitute the \emph{two}-predictor model for $Y$; \mv{3} split the covariance; divide by $\Var X_1$; name the leftover ratio $\beta_{1,2}$.}
\end{methodbox}
\begin{align*}
\Cov(X_1,Y)&=\Cov(X_1,\ \beta_0+\beta_1X_1+\beta_2X_2+\epsilon) \why{substitute the two-predictor model}\\
&=\beta_1\Cov(X_1,X_1)+\beta_2\Cov(X_1,X_2)+\Cov(X_1,\epsilon) \why{bilinearity; $\beta_0$ drops out}\\
&=\beta_1\Var(X_1)+\beta_2\Cov(X_1,X_2). \why{$\Cov(X_1,\epsilon)=0$}
\end{align*}
Divide by $\Var(X_1)$:
\[
\beta_1'=\frac{\Cov(X_1,Y)}{\Var(X_1)}=\beta_1+\beta_2\,\frac{\Cov(X_1,X_2)}{\Var(X_1)} .
\]
The last fraction has the same form as Section~\ref{sec:beta1} with $X_2$ in the role of $Y$: it is the slope when $X_2$ is regressed on $X_1$.
\begin{keybox}[title={Omitted-variable bias}]
\[
\beta_1'=\beta_1+\beta_2\,\beta_{1,2},\qquad \beta_{1,2}=\frac{\Cov(X_1,X_2)}{\Var(X_1)}\ \ \text{(slope of $X_2$ on $X_1$)}.
\]
The bias $\beta_2\beta_{1,2}$ is $0$ only if $\beta_2=0$ (no effect) or $\beta_{1,2}=0$ (uncorrelated, as in a randomized trial).
\end{keybox}
(The notes assume mean-zero predictors and no intercept ``for simplicity''. As the steps above show, that is not needed: constants
always drop out of a covariance.)

\textbf{Method 2: conditional means (when $X_1$ is binary).} This is how Lecture 12 explained the test-score example. Condition on $X_1$ \mv{5}
and take $\E$ \mv{2}:
\[
\E[Y\given X_1]=\beta_0+\beta_1X_1+\beta_2\E[X_2\given X_1]+\underbrace{\E[\epsilon\given X_1]}_{0}.
\]
Subtract the $X_1=0$ version from the $X_1=1$ version (as in Section~\ref{sec:binary}):
\[
\underbrace{\E[Y\given X_1{=}1]-\E[Y\given X_1{=}0]}_{\beta_1'}=\beta_1+\beta_2\underbrace{\big(\E[X_2\given X_1{=}1]-\E[X_2\given X_1{=}0]\big)}_{\beta_{1,2}} .
\]
Same formula. Check with the numbers: $11.8\approx5.95+0.56\times10.3$.

\begin{tipbox}[title={Trick from the crime example: two OVB slopes give the predictor correlation}]
$\beta_{1,2}\,\beta_{2,1}=\dfrac{\Cov(X_1,X_2)}{\Var X_1}\cdot\dfrac{\Cov(X_1,X_2)}{\Var X_2}=\left(\dfrac{\Cov(X_1,X_2)}{\sigma_{X_1}\sigma_{X_2}}\right)^2=\rho_{1,2}^2 .$
So if you solve OVB both ways for the two slopes, $\rho_{1,2}=\pm\sqrt{\beta_{1,2}\beta_{2,1}}$, with the sign of the slopes.
\end{tipbox}

\subsection{(extra) Omitted-variable bias with $K$ predictors}\label{sec:ovbk}
\begin{methodbox}
\goal{$\beta_1'$ when $X_2,\dots,X_K$ are all left out.}
\start{$\beta_1'=\Cov(X_1,Y)/\Var(X_1)$.}
\moves{identical to Section~\ref{sec:ovb}: \mv{1}\,\mv{3}, now with a sum of left-out terms.}
\end{methodbox}
\begin{align*}
\Cov(X_1,Y)&=\Cov\Big(X_1,\ \beta_0+\beta_1X_1+\sum_{j=2}^K\beta_jX_j+\epsilon\Big) \why{substitute the model}\\
&=\beta_1\Var(X_1)+\sum_{j=2}^K\beta_j\Cov(X_1,X_j)+0 \why{bilinearity, one term at a time}
\end{align*}
\begin{keybox}
\[\beta_1'=\beta_1+\sum_{j=2}^K\beta_j\,\beta_{1,j},\qquad\beta_{1,j}=\frac{\Cov(X_1,X_j)}{\Var(X_1)} .\]
Picture: every path from $X_1$ to $Y$ adds up: the direct arrow $\beta_1$, plus each detour $X_1\to X_j\to Y$ worth $\beta_{1,j}\beta_j$.
\end{keybox}

\subsection{(extra) The covariance system: solving for $\beta_1$ and $\beta_2$}\label{sec:system}
\begin{methodbox}
\goal{the two-predictor coefficients in terms of covariances.}
\start{the two-predictor model.}
\moves{\mv{3} \emph{twice}: take $\Cov(X_1,\cdot)$ and $\Cov(X_2,\cdot)$ of the model. That gives two equations in two unknowns; eliminate $\beta_2$.}
\end{methodbox}
Write $c=\Cov(X_1,X_2)$. The first line was done in Section~\ref{sec:ovb}; the second is the same move with $X_2$:
\begin{align*}
\Cov(X_1,Y)&=\beta_1\sigma_{X_1}^2+\beta_2\,c \tag{i}\\
\Cov(X_2,Y)&=\beta_1\,c+\beta_2\sigma_{X_2}^2 \tag{ii}
\end{align*}
\textbf{Eliminate $\beta_2$:} multiply (i) by $\sigma_{X_2}^2$ and (ii) by $c$, then subtract:
\[
\sigma_{X_2}^2\Cov(X_1,Y)-c\,\Cov(X_2,Y)=\beta_1\big(\sigma_{X_1}^2\sigma_{X_2}^2-c^2\big).
\]
\begin{keybox}
\[
\beta_1=\frac{\Cov(X_1,Y)\,\sigma_{X_2}^2-\Cov(X_2,Y)\,\Cov(X_1,X_2)}{\sigma_{X_1}^2\sigma_{X_2}^2-\Cov(X_1,X_2)^2}
\ \overset{\text{standardized}}{=}\ \frac{\rho_1-\rho_{1,2}\,\rho_2}{1-\rho_{1,2}^2},
\]
where $\rho_k=\corr(X_k,Y)$ and $\rho_{1,2}=\corr(X_1,X_2)$.
\end{keybox}
For the standardized version, put all of $X_1,X_2,Y$ in SD units \mv{6}: then every $\sigma=1$, $\Cov(X_k,Y)=\rho_k$, and $c=\rho_{1,2}$.
\textbf{Checks:} (1) dividing (i) by $\sigma_{X_1}^2$ is exactly OVB; (2) if $\rho_{1,2}=0$ then $\beta_1=\rho_1$, the single-predictor slope.

\textbf{Matrix form (as in the notes).} (i)--(ii) say $\big(\Cov(X_1,Y),\Cov(X_2,Y)\big)^T=\Sigma\,(\beta_1,\beta_2)^T$ with
$\Sigma=\begin{bmatrix}\sigma_{X_1}^2&c\\c&\sigma_{X_2}^2\end{bmatrix}$, so $(\beta_1,\beta_2)^T=\Sigma^{-1}\big(\Cov(X_1,Y),\Cov(X_2,Y)\big)^T$
using $\begin{bmatrix}a&b\\b&d\end{bmatrix}^{-1}=\frac{1}{ad-b^2}\begin{bmatrix}d&-b\\-b&a\end{bmatrix}$. It gives the same $\beta_1$.

\subsection{(extra) Correlated predictors inflate the standard error}\label{sec:inflate}
\begin{methodbox}
\goal{$\se(\hat\beta_1)$ when $X_2$ is also in the model.}
\start{the result stated in the notes (Unit 4.3--4.4): the variances and covariance of $(\hat\beta_1,\hat\beta_2)$ are $\frac{\sigma_\epsilon^2}{N}\Sigma^{-1}$.}
\moves{invert $\Sigma$ (2$\times$2 rule); read off the diagonal for variances and the off-diagonal for the covariance; simplify with
$c=\rho_{1,2}\sigma_{X_1}\sigma_{X_2}$.}
\end{methodbox}
Since $c=\rho_{1,2}\sigma_{X_1}\sigma_{X_2}$, the determinant is $\sigma_{X_1}^2\sigma_{X_2}^2-c^2=\sigma_{X_1}^2\sigma_{X_2}^2(1-\rho_{1,2}^2)$. The top-left entry of $\Sigma^{-1}$ is
\[
\frac{\sigma_{X_2}^2}{\sigma_{X_1}^2\sigma_{X_2}^2(1-\rho_{1,2}^2)}=\frac{1}{\sigma_{X_1}^2(1-\rho_{1,2}^2)}
\quad\Longrightarrow\quad
\Var(\hat\beta_1)=\frac{\sigma_\epsilon^2}{N\sigma_{X_1}^2(1-\rho_{1,2}^2)} .
\]
\begin{keybox}
\[
\se(\hat\beta_1)=\frac{\sigma_\epsilon}{\sigma_{X_1}\sqrt{N(1-\rho_{1,2}^2)}}=\underbrace{\frac{\sigma_\epsilon}{\sigma_{X_1}\sqrt N}}_{\text{one predictor}}\cdot\frac{1}{\sqrt{1-\rho_{1,2}^2}},
\qquad \corr(\hat\beta_1,\hat\beta_2)=-\rho_{1,2}.
\]
\end{keybox}
\textbf{Correlation of the estimates} (standardized predictors): $\Sigma^{-1}=\frac{1}{1-\rho_{1,2}^2}\begin{bmatrix}1&-\rho_{1,2}\\-\rho_{1,2}&1\end{bmatrix}$,
so $\corr=\dfrac{\Cov}{\se\cdot\se}=\dfrac{-\rho_{1,2}/(1-\rho_{1,2}^2)}{1/(1-\rho_{1,2}^2)}=-\rho_{1,2}$.

\textbf{Method 2 (intuition you can rebuild from the sheet).} Only the part of $X_1$ that $X_2$ cannot predict helps separate $\beta_1$ from
$\beta_2$. Apply the variance-explained formula (Section~\ref{sec:rho2}) to $X_1$ regressed on $X_2$: the leftover variance is
$\Var(X_1\given X_2)=(1-\rho_{1,2}^2)\sigma_{X_1}^2$. Put that leftover SD in place of $\sigma_{X_1}$ in $\se=\sigma_\epsilon/(\sigma_{X_1}\sqrt N)$ and you get
the boxed formula.

%==========================================================
\clearpage
\section{Estimators and their sample distributions}\label{sec:est}
%==========================================================
\subsection{Definitions (nothing to derive)}
$Y_1,\dots,Y_N$ iid; an estimator $\hat\theta=\hat\theta(Y_1,\dots,Y_N)$ is any function of the data. Over replicate datasets it varies; its
distribution over replicates is the \textbf{sample distribution} and its standard deviation is $\se(\hat\theta)=\sqrt{\Var(\hat\theta)}$.
$\bias(\hat\theta)=\E[\hat\theta]-\theta$. These are definitions: learn them, do not derive them.

\subsection{The sample mean: $\Ybar\approx\Norm(\mu,\sigma^2/N)$, $\se(\Ybar)=\sigma/\sqrt N$}\label{sec:ybar}
\begin{methodbox}
\goal{the sample distribution of $\Ybar$ when $\E Y_i=\mu$, $\Var Y_i=\sigma^2$.}
\start{$\Ybar=\frac1N\sum_iY_i$.}
\moves{the general recipe: \mv{2} center, \mv{4} spread (independence: variances add), \mv{7} shape.}
\end{methodbox}
\textbf{Center.}
\[
\E[\Ybar]=\E\Big[\frac1N\sum_iY_i\Big]=\frac1N\sum_i\E[Y_i]=\frac1N\cdot N\mu=\mu .
\]
\textbf{Spread.}
\begin{align*}
\Var(\Ybar)&=\Var\Big(\frac1N\sum_iY_i\Big)=\frac{1}{N^2}\Var\Big(\sum_iY_i\Big) \why{$\Var(aX)=a^2\Var X$ with $a=1/N$}\\
&=\frac{1}{N^2}\sum_i\Var(Y_i)=\frac{1}{N^2}\cdot N\sigma^2=\frac{\sigma^2}{N} \why{independent: variances add}
\end{align*}
\textbf{Shape.} If the $Y_i$ are Normal, $\sum_iY_i$ is exactly Normal (a sum of independent Normals). Otherwise the CLT says
$\Ybar$ is approximately Normal for large $N$.
\begin{keybox}
\[\Ybar\approx\Norm\Big(\mu,\ \frac{\sigma^2}{N}\Big),\qquad\se(\Ybar)=\frac{\sigma}{\sqrt N}.\]
Unbiased ($\E\Ybar=\mu$) and consistent ($\se\to0$).
\end{keybox}
\textbf{Bernoulli version.} If $Y_i\sim\Bern(q)$ then $\mu=q$ and $\sigma^2=q(1-q)$, so $\hat q=\Ybar\approx\Norm\big(q,\ q(1-q)/N\big)$.

\textbf{Common slips:} the $1/N$ comes out of the variance \emph{squared} ($1/N^2$); $\se$ has $\sqrt N$, the variance has $N$.

\subsection{Bias--variance decomposition: $\MSE=\se^2+\bias^2$}\label{sec:mse}
\begin{methodbox}
\goal{split $\MSE_{\hat\theta}=\E[(\hat\theta-\theta)^2]$ into variance and bias.}
\start{either the variance formula $\Var(W)=\E[W^2]-(\E W)^2$ (method A, the notes' proof) or the square itself (method B).}
\moves{A: apply the formula to $W=\hat\theta-\theta$. \ B: \mv{10} add and subtract $\E[\hat\theta]$, expand, show the cross term is $0$.}
\end{methodbox}
\textbf{Method A.} Let $W=\hat\theta-\theta$. Since $\theta$ is a fixed number, $\Var(W)=\Var(\hat\theta)$ and $\E W=\E\hat\theta-\theta$.
\begin{align*}
\Var(W)&=\E[W^2]-(\E W)^2 \why{variance formula}\\
\Var(\hat\theta)&=\E\big[(\hat\theta-\theta)^2\big]-\big(\E[\hat\theta]-\theta\big)^2 \why{shifting by $\theta$ does not change variance}\\
\se(\hat\theta)^2&=\MSE_{\hat\theta}-\bias(\hat\theta)^2 . \why{definitions}
\end{align*}
\textbf{Method B.} Write $m=\E[\hat\theta]$ (a constant).
\begin{align*}
\E\big[(\hat\theta-\theta)^2\big]&=\E\big[\big((\hat\theta-m)+(m-\theta)\big)^2\big] \why{add and subtract $m$}\\
&=\E\big[(\hat\theta-m)^2\big]+2(m-\theta)\,\E[\hat\theta-m]+(m-\theta)^2 \why{expand; $m-\theta$ is a constant}\\
&=\Var(\hat\theta)+0+(m-\theta)^2 \why{$\E[\hat\theta-m]=m-m=0$}
\end{align*}
\begin{keybox}
\[\MSE_{\hat\theta}=\se(\hat\theta)^2+\bias(\hat\theta)^2 .\]
\end{keybox}

\subsection{(extra) Bessel's correction: why divide by $N-1$}\label{sec:bessel}
\begin{methodbox}
\goal{$\E\big[\sum_i(Y_i-\Ybar)^2\big]$.}
\start{the sum of squares $\sum_i(Y_i-\Ybar)^2$.}
\moves{\mv{10} add and subtract $\mu$ inside the square; expand; use $\sum_i(Y_i-\mu)=N(\Ybar-\mu)$; take $\E$ using
$\E[(Y_i-\mu)^2]=\sigma^2$ and $\E[(\Ybar-\mu)^2]=\Var\Ybar=\sigma^2/N$.}
\end{methodbox}
\begin{align*}
\sum_i(Y_i-\Ybar)^2&=\sum_i\big[(Y_i-\mu)-(\Ybar-\mu)\big]^2 \why{add and subtract $\mu$}\\
&=\sum_i(Y_i-\mu)^2-2(\Ybar-\mu)\sum_i(Y_i-\mu)+N(\Ybar-\mu)^2 \why{expand; $\Ybar-\mu$ same for all $i$}\\
&=\sum_i(Y_i-\mu)^2-2N(\Ybar-\mu)^2+N(\Ybar-\mu)^2 \why{$\sum_i(Y_i-\mu)=N(\Ybar-\mu)$}\\
&=\sum_i(Y_i-\mu)^2-N(\Ybar-\mu)^2 .
\end{align*}
Take expectations: $\E\big[\sum_i(Y_i-\Ybar)^2\big]=N\sigma^2-N\cdot\frac{\sigma^2}{N}=(N-1)\sigma^2$.
\begin{keybox}
\[\hat\sigma_0^2=\tfrac1N\textstyle\sum_i(Y_i-\Ybar)^2\text{ has }\E[\hat\sigma_0^2]=\tfrac{N-1}{N}\sigma^2;\qquad
\hat\sigma^2=\tfrac{1}{N-1}\textstyle\sum_i(Y_i-\Ybar)^2\text{ is unbiased.}\]
\end{keybox}
The same idea explains $\hat\sigma_\epsilon^2=\SSR/(N-2)$: two parameters were fitted from the data, so divide by $N-2$.

%==========================================================
\clearpage
\section{Confidence intervals and hypothesis tests}\label{sec:citest}
%==========================================================
\subsection{The 95\% confidence interval: $\hat\theta\pm1.96\,\se(\hat\theta)$}\label{sec:ci}
\begin{methodbox}
\goal{an interval that contains $\theta$ in 95\% of replicates.}
\start{the sample distribution $\hat\theta\approx\Norm(\theta,\se^2)$. This needs $\hat\theta$ \textbf{unbiased}, so the center is $\theta$.}
\moves{\mv{6} standardize; \mv{9} 95\% of a standard Normal lies in $(-1.96,1.96)$; rearrange the inequality to isolate $\theta$.}
\end{methodbox}
\begin{align*}
&\frac{\hat\theta-\theta}{\se(\hat\theta)}\approx\Norm(0,1) \why{standardize; uses $\E\hat\theta=\theta$}\\
&\P\Big(-1.96<\frac{\hat\theta-\theta}{\se(\hat\theta)}<1.96\Big)\approx0.95 \why{$\P(|Z|<1.96)=0.95$}\\
&\P\big(-1.96\,\se<\hat\theta-\theta<1.96\,\se\big)\approx0.95 \why{multiply through by $\se>0$}\\
&\P\big(\hat\theta-1.96\,\se<\theta<\hat\theta+1.96\,\se\big)\approx0.95 \why{rearrange both inequalities for $\theta$}
\end{align*}
\begin{keybox}
\[\text{95\% CI}=\big[\hat\theta-1.96\,\se(\hat\theta),\ \hat\theta+1.96\,\se(\hat\theta)\big]\qquad\text{(when $\hat\theta$ is unbiased).}\]
The random thing is the interval, not $\theta$: over replicates, 95\% of such intervals contain $\theta$.
\end{keybox}
\textbf{Why ``unbiased'' matters.} If $\E\hat\theta\ne\theta$, the bell curve is centered in the wrong place and the step ``standardize'' fails:
the interval covers $\theta$ less than 95\% of the time.

\textbf{(extra) How much data for a target width.} If $\se=\sigma/\sqrt N$ and you want half-width at most $h$:
$1.96\,\sigma/\sqrt N\le h\iff\sqrt N\ge1.96\,\sigma/h\iff N\ge(1.96\,\sigma/h)^2$.

\subsection{The test statistic and $p$-value}\label{sec:test}
\begin{methodbox}
\goal{how surprising $\hat\beta_1$ would be if $H_0:\beta_1=0$ were true.}
\start{the sample distribution $\hat\beta_1\approx\Norm(\beta_1,\se^2)$ (Section~\ref{sec:varb1}).}
\moves{set $\beta_1=0$ (the null); \mv{6} standardize by $\se$; \mv{9} two tails of a standard Normal.}
\end{methodbox}
Under $H_0$, a replicate estimate $\hat\beta_1^*\approx\Norm(0,\se(\hat\beta_1)^2)$. Dividing by $\se$ \mv{6}:
\[
\hat T^*=\frac{\hat\beta_1^*}{\se(\hat\beta_1)}\approx\Norm(0,1),\qquad\text{and for our data}\qquad \hat T=\frac{\hat\beta_1}{\se(\hat\beta_1)} .
\]
The $p$-value is the chance a null replicate is at least as extreme as ours. Split it into two tails and use symmetry of $\Norm(0,1)$:
\begin{align*}
p_v=\P\big(|\hat T^*|>|\hat T|\big)&=\P\big(\hat T^*>|\hat T|\big)+\P\big(\hat T^*<-|\hat T|\big) \why{two tails}\\
&=2\big(1-\Phi(|\hat T|)\big) \why{both tails have the same area}
\end{align*}
\begin{keybox}
\[H_0:\beta_1=0,\qquad \hat T=\frac{\hat\beta_1}{\se(\hat\beta_1)},\qquad p_v=\P\big(|\hat T^*|>|\hat T|\big)=2\big(1-\Phi(|\hat T|)\big).\]
\end{keybox}
\textbf{By hand:} near $0$ the standard Normal density is about $0.4$, so $\P(|Z|<a)\approx0.8a$ for small $a$; then $p_v\approx1-0.8|\hat T|$.

\subsection{$p_v<0.05\iff0\notin$ 95\% CI}\label{sec:ci-test}
\begin{methodbox}
\goal{show the test and the CI always agree.}
\start{$\P(|Z|>1.96)=0.05$, and the larger $|\hat T|$ is, the smaller $p_v$ is.}
\moves{a chain of equivalent statements: $p_v$ $\to$ $|\hat T|$ $\to$ $|\hat\beta_1|$ $\to$ the CI endpoints.}
\end{methodbox}
\begin{align*}
p_v<0.05
&\iff|\hat T|>1.96 \why{$p_v=0.05$ at $|\hat T|=1.96$; $p_v$ shrinks as $|\hat T|$ grows}\\
&\iff|\hat\beta_1|>1.96\,\se(\hat\beta_1) \why{multiply by $\se>0$}\\
&\iff\hat\beta_1-1.96\,\se>0\ \text{ or }\ \hat\beta_1+1.96\,\se<0 \why{$\hat\beta_1>1.96\,\se$ or $\hat\beta_1<-1.96\,\se$}\\
&\iff\text{CI entirely above 0 or entirely below 0} \why{these are the CI endpoints}\\
&\iff0\notin\text{95\% CI}.
\end{align*}

%==========================================================
\clearpage
\section{Least squares and the sample distribution of $\hat\beta_1$}\label{sec:ls}
%==========================================================
Data $(X_1,Y_1),\dots,(X_N,Y_N)$ from the one-predictor model. (Here $X_i$ means the $i$-th data point of the single predictor.)

\subsection{The estimators $\hat\beta_1$ and $\hat\beta_0$}\label{sec:lsest}
\begin{methodbox}
\goal{estimates of $\beta_1,\beta_0$ from data.}
\start{the population formulas $\beta_1=\Cov(X,Y)/\sigma_X^2$ and $\beta_0=\E Y-\beta_1\E X$ (Sections~\ref{sec:beta1}--\ref{sec:beta0}).}
\moves{\mv{8} plug in: replace every $\E[\cdot]$ by $\frac1N\sum_i$ and $\E X,\E Y$ by $\Xbar,\Ybar$. The $\frac1N$'s cancel.}
\end{methodbox}
Write $\Cov(X,Y)=\E[(X-\E X)(Y-\E Y)]$ and $\sigma_X^2=\E[(X-\E X)^2]$, then plug in:
\[
\hat\beta_1=\frac{\frac1N\sum_i(X_i-\Xbar)(Y_i-\Ybar)}{\frac1N\sum_i(X_i-\Xbar)^2}=\frac{\sum_i(X_i-\Xbar)(Y_i-\Ybar)}{\sum_i(X_i-\Xbar)^2},
\qquad \hat\beta_0=\Ybar-\hat\beta_1\Xbar .
\]
\textbf{Method 2 (calculus; the ``least squares'' name).} Minimize $f(b_0,b_1)=\sum_i(Y_i-b_0-b_1X_i)^2$: set the derivative in each of $b_0,b_1$ to $0$.
\begin{align*}
\text{in }b_0:\quad&\textstyle\sum_i-2(Y_i-b_0-b_1X_i)=0\ \Longrightarrow\ N\Ybar-Nb_0-b_1N\Xbar=0\ \Longrightarrow\ b_0=\Ybar-b_1\Xbar,\\
\text{in }b_1:\quad&\textstyle\sum_i-2X_i(Y_i-b_0-b_1X_i)=0 .
\end{align*}
Substitute $b_0$ into the second equation: $\sum_iX_i\big[(Y_i-\Ybar)-b_1(X_i-\Xbar)\big]=0$, so $b_1=\frac{\sum_iX_i(Y_i-\Ybar)}{\sum_iX_i(X_i-\Xbar)}$.
Then the sum trick \mv{11}: subtracting $\Xbar\sum_i(Y_i-\Ybar)=0$ on top and $\Xbar\sum_i(X_i-\Xbar)=0$ on the bottom turns $X_i$ into
$(X_i-\Xbar)$, giving the same $\hat\beta_1$.
\begin{keybox}
\[\hat\beta_1=\frac{\sum_i(X_i-\Xbar)(Y_i-\Ybar)}{\sum_i(X_i-\Xbar)^2},\qquad\hat\beta_0=\Ybar-\hat\beta_1\Xbar .\]
\end{keybox}

\subsection{(extra) $\hat\sigma_\epsilon^2$ and $R^2$}\label{sec:r2}
\begin{methodbox}
\goal{estimates of $\sigma_\epsilon^2$ and $\rho^2$.}
\start{$\sigma_\epsilon^2=\Var(Y\given X)$ and $\rho^2=1-\Var(Y\given X)/\Var(Y)$.}
\moves{\mv{8} plug in: the spread around the line is estimated by the residuals; the total spread by deviations from $\Ybar$.}
\end{methodbox}
Residuals $r_i=Y_i-\hat\beta_0-\hat\beta_1X_i$, $\SSR=\sum_ir_i^2$, $\SST=\sum_i(Y_i-\Ybar)^2$. Then $\Var(Y\given X)\approx\SSR/N$ and
$\Var(Y)\approx\SST/N$, and the $N$'s cancel in the ratio:
\begin{keybox}
\[\hat\sigma_\epsilon^2=\frac{\SSR}{N-2}\ \Big(\approx\frac{\SSR}{N}\Big),\qquad R^2=1-\frac{\SSR}{\SST}\approx\rho^2 .\]
\end{keybox}
(The $N-2$ instead of $N$ is Bessel's correction again, Section~\ref{sec:bessel}.)

\subsection{$\hat\beta_1$ is unbiased: $\E[\hat\beta_1]=\beta_1$}\label{sec:unbiased}
\begin{methodbox}
\goal{$\E[\hat\beta_1\given X_1,\dots,X_N]$.}
\start{the formula for $\hat\beta_1$, treating the $X$'s as fixed numbers \mv{5}.}
\moves{\mv{2} linearity moves $\E$ onto the $Y$'s; \mv{1} $\E[Y_i\given X_i]=\beta_0+\beta_1X_i$; cancel.}
\end{methodbox}
Given the $X$'s, the weights $(X_i-\Xbar)$ and the denominator $S_{xx}=\sum_i(X_i-\Xbar)^2$ are constants.
\begin{align*}
\E[\hat\beta_1\given X]&=\frac{\sum_i(X_i-\Xbar)\,\E[Y_i-\Ybar\given X]}{S_{xx}} \why{linearity; constants come out}\\
\E[Y_i-\Ybar\given X]&=(\beta_0+\beta_1X_i)-(\beta_0+\beta_1\Xbar)=\beta_1(X_i-\Xbar) \why{model, averaged over $i$ for $\Ybar$}\\
\E[\hat\beta_1\given X]&=\frac{\sum_i(X_i-\Xbar)\,\beta_1(X_i-\Xbar)}{S_{xx}}=\beta_1\frac{S_{xx}}{S_{xx}}=\beta_1 . \why{substitute}
\end{align*}
\begin{keybox}
\[\E[\hat\beta_1]=\beta_1\qquad\text{(conditional on the $X$'s).}\]
The key ingredient is that the noise is independent of \emph{all} the predictors (this fails for AR models: the Hurwicz bias).
\end{keybox}

\subsection{The spread: $\se(\hat\beta_1)\approx\sigma_\epsilon/(\sigma_{X}\sqrt N)$}\label{sec:varb1}
\begin{methodbox}
\goal{$\Var(\hat\beta_1\given X)$, then the full sample distribution.}
\start{$\hat\beta_1$ written as a weighted sum of the $Y_i$.}
\moves{\mv{11} drop $\Ybar$ from the numerator; \mv{4} independent $Y_i$'s, each with $\Var(Y_i\given X_i)=\sigma_\epsilon^2$: variances add with weights squared;
replace $S_{xx}$ by $N\sigma_X^2$; shape from ``a weighted sum of Normals is Normal''.}
\end{methodbox}
\textbf{Step 1: weighted sum.} $\sum_i(X_i-\Xbar)\Ybar=\Ybar\sum_i(X_i-\Xbar)=0$ \mv{11}, so
\[
\hat\beta_1=\frac{\sum_i(X_i-\Xbar)Y_i}{S_{xx}}=\sum_iw_iY_i,\qquad w_i=\frac{X_i-\Xbar}{S_{xx}} .
\]
\textbf{Step 2: variance.}
\begin{align*}
\Var(\hat\beta_1\given X)&=\sum_iw_i^2\Var(Y_i\given X_i) \why{independent $Y_i$; $\Var(w_iY_i)=w_i^2\Var Y_i$}\\
&=\frac{\sum_i(X_i-\Xbar)^2}{S_{xx}^2}\,\sigma_\epsilon^2=\frac{S_{xx}}{S_{xx}^2}\,\sigma_\epsilon^2=\frac{\sigma_\epsilon^2}{S_{xx}} \why{$\Var(Y_i\given X_i)=\sigma_\epsilon^2$}\\
&\approx\frac{\sigma_\epsilon^2}{N\sigma_X^2} \why{$S_{xx}/N$ is the sample variance of $X$, $\approx\sigma_X^2$}
\end{align*}
\textbf{Step 3: shape.} Given the $X$'s, each $Y_i$ is Normal, so the weighted sum $\hat\beta_1$ is Normal.
\begin{keybox}
\[\E[\hat\beta_1]=\beta_1,\qquad \se(\hat\beta_1)\approx\frac{\sigma_\epsilon}{\sigma_{X}\sqrt N},\qquad
\hat\beta_1\approx\Norm\Big(\beta_1,\ \frac{\sigma_\epsilon^2}{N\sigma_{X}^2}\Big).\]
More precise with less noise, more spread-out $X$'s, or more data.
\end{keybox}
\textbf{Compare with $\Ybar$.} The recipe is identical: $\Ybar$ has weights $1/N$ (giving $\sigma^2/N$); $\hat\beta_1$ has weights $w_i$ (giving $\sigma_\epsilon^2/S_{xx}$).

\subsection{(extra) Check: a binary predictor gives a difference of means}\label{sec:binarycheck}
\begin{methodbox}
\goal{the same $\se$ by a second route when $X\sim\Bern(p)$.}
\start{$\hat\beta_1=\overline{Y}_{X=1}-\overline{Y}_{X=0}$ (Section~\ref{sec:binary}).}
\moves{each group mean is $\Ybar$ for that group \mv{7}; the groups are independent, so variances add \mv{4} (even though we subtract).}
\end{methodbox}
With $Np$ points in group 1 and $N(1-p)$ in group 0,
\begin{align*}
\Var(\hat\beta_1)&=\frac{\sigma_\epsilon^2}{Np}+\frac{\sigma_\epsilon^2}{N(1-p)} \why{$\Var(A-B)=\Var A+\Var B$ for independent $A,B$}\\
&=\sigma_\epsilon^2\,\frac{(1-p)+p}{Np(1-p)}=\frac{\sigma_\epsilon^2}{N\,p(1-p)}=\frac{\sigma_\epsilon^2}{N\sigma_X^2} \why{$\sigma_X^2=p(1-p)$ for a Bernoulli}
\end{align*}
This matches Section~\ref{sec:varb1}. With equal groups ($p=\frac12$) it is $4\sigma_\epsilon^2/N$, the clinical-trial result.

\subsection{(extra) $\hat T^2=(N-2)\,R^2/(1-R^2)$}\label{sec:tr2}
\begin{methodbox}
\goal{connect the test statistic to $R^2$.}
\start{$r_i=(Y_i-\Ybar)-\hat\beta_1(X_i-\Xbar)$ (from $\hat\beta_0=\Ybar-\hat\beta_1\Xbar$), and $\se(\hat\beta_1)^2=\hat\sigma_\epsilon^2/S_{xx}$.}
\moves{square and sum $r_i$ to get $\SSR=\SST-\hat\beta_1^2S_{xx}$; substitute into $\hat T^2$; divide top and bottom by $\SST$.}
\end{methodbox}
Using $\sum_i(X_i-\Xbar)(Y_i-\Ybar)=\hat\beta_1S_{xx}$ (the definition of $\hat\beta_1$):
\begin{align*}
\SSR&=\SST-2\hat\beta_1\sum_i(X_i-\Xbar)(Y_i-\Ybar)+\hat\beta_1^2S_{xx} \why{expand $\sum_ir_i^2$}\\
&=\SST-2\hat\beta_1^2S_{xx}+\hat\beta_1^2S_{xx}=\SST-\hat\beta_1^2S_{xx}. \why{substitute}
\end{align*}
So $\hat\beta_1^2S_{xx}=\SST-\SSR$. With $\se^2=\SSR/((N-2)S_{xx})$:
\begin{align*}
\hat T^2=\frac{\hat\beta_1^2}{\se^2}&=\frac{(N-2)\hat\beta_1^2S_{xx}}{\SSR}=\frac{(N-2)(\SST-\SSR)}{\SSR} \why{substitute}\\
&=(N-2)\frac{1-\SSR/\SST}{\SSR/\SST}=(N-2)\frac{R^2}{1-R^2}. \why{divide top and bottom by $\SST$}
\end{align*}
\textbf{Use:} a tiny $p_v$ with a small $R^2$ is possible (large $N$); a large $p_v$ with $R^2$ near 1 is not.

%==========================================================
\clearpage
\section{Self-test: cover the right two columns}\label{sec:selftest}
%==========================================================
For each formula, say out loud what you start from and the key move, then check. If you can do every row, you can derive the whole sheet.
{\small
\renewcommand{\arraystretch}{1.28}
\begin{longtable}{@{}>{\raggedright\arraybackslash}p{4.6cm}>{\raggedright\arraybackslash}p{4.4cm}>{\raggedright\arraybackslash}p{6.6cm}@{}}
\toprule
\textbf{Result} & \textbf{Start from} & \textbf{Key move(s)}\\
\midrule
\endhead
$\beta_1=\Cov(X_1,Y)/\sigma_{X_1}^2$ & model & $\Cov(X_1,\cdot)$ of both sides; $\beta_0$ and $\epsilon$ drop out\\
$\beta_0=\E Y-\beta_1\E X_1$ & model & $\E$ of both sides; $\E\epsilon=0$\\
binary $X$: $\beta_1=$ difference of group means & model given $X=1$ and $X=0$ & condition, take $\E$, subtract\\
$\sigma_Y^2=\beta_1^2\sigma_{X_1}^2+\sigma_\epsilon^2$ & model & $\Var$ of both sides; independent pieces add; coefficient squared\\
$\Var(Y\given X_1)=\sigma_\epsilon^2$ & model & given $X_1$, $\beta_0+\beta_1X_1$ is a constant\\
$\rho=\beta_1\sigma_{X_1}/\sigma_Y$ & $\rho=\Cov/(\sigma_X\sigma_Y)$, or standardized model & substitute $\Cov=\beta_1\sigma_X^2$; or standardize and read the slope\\
$\rho^2=1-\Var(Y\given X)/\Var Y$ & variance equation & divide by $\sigma_Y^2$; first term is $\rho^2$\\
$\rho^2$ never decreases & $\Var(Y\given X_1)$ in the 2-predictor model & $=\beta_2^2\Var(X_2\given X_1)+\sigma_\epsilon^2\ge\sigma_\epsilon^2$\\
OVB $\beta_1'=\beta_1+\beta_2\beta_{1,2}$ & $\beta_1'=\Cov(X_1,Y)/\Var X_1$ & substitute the \emph{2-predictor} model; bilinearity; divide\\
OVB, $K$ predictors & same & same, one term per left-out $X_j$\\
$\beta_{1,2}\beta_{2,1}=\rho_{1,2}^2$ & definitions of the slopes & multiply; $\Cov^2/(\Var\Var)$\\
$\beta_1=\frac{\rho_1-\rho_{1,2}\rho_2}{1-\rho_{1,2}^2}$ & model & $\Cov(X_1,\cdot)$ \emph{and} $\Cov(X_2,\cdot)$; 2 equations; eliminate $\beta_2$\\
$\se$ inflation $1/\sqrt{1-\rho_{1,2}^2}$ & $\frac{\sigma_\epsilon^2}{N}\Sigma^{-1}$ (given) & 2$\times$2 inverse; determinant $=\sigma_1^2\sigma_2^2(1-\rho_{1,2}^2)$\\
$\Ybar\approx\Norm(\mu,\sigma^2/N)$ & $\Ybar=\frac1N\sum Y_i$ & $\E$: linearity; $\Var$: $\frac1{N^2}\cdot N\sigma^2$; shape: CLT\\
$\MSE=\se^2+\bias^2$ & $\Var W=\E W^2-(\E W)^2$ & $W=\hat\theta-\theta$; or add and subtract $\E\hat\theta$\\
Bessel ($N-1$) & $\sum(Y_i-\Ybar)^2$ & add and subtract $\mu$; $=\sum(Y_i-\mu)^2-N(\Ybar-\mu)^2$\\
95\% CI $\hat\theta\pm1.96\,\se$ & $\hat\theta\approx\Norm(\theta,\se^2)$ (unbiased) & standardize; $\P(|Z|<1.96)=.95$; solve for $\theta$\\
$\hat T$, $p_v=2(1-\Phi(|\hat T|))$ & $\hat\beta_1^*\approx\Norm(0,\se^2)$ under $H_0$ & divide by $\se$; two equal tails\\
$p_v<.05\iff0\notin$ CI & $\P(|Z|>1.96)=.05$ & $|\hat T|>1.96\iff|\hat\beta_1|>1.96\,\se\iff$ CI misses 0\\
$\hat\beta_1,\hat\beta_0$ & $\beta_1=\Cov/\sigma_X^2$, $\beta_0=\E Y-\beta_1\E X$ & plug in sample averages; or minimize SSR\\
$\hat\sigma_\epsilon^2$, $R^2$ & $\sigma_\epsilon^2=\Var(Y\given X)$, $\rho^2$ & plug in $\SSR/N$ and $\SST/N$\\
$\E\hat\beta_1=\beta_1$ & $\hat\beta_1$ formula, $X$'s fixed & linearity; $\E[Y_i-\Ybar\given X]=\beta_1(X_i-\Xbar)$\\
$\se(\hat\beta_1)\approx\sigma_\epsilon/(\sigma_X\sqrt N)$ & $\hat\beta_1=\sum w_iY_i$ & sum trick drops $\Ybar$; $\Var=\sum w_i^2\sigma_\epsilon^2=\sigma_\epsilon^2/S_{xx}$\\
binary-$X$ check & difference of group means & independent groups: variances add; $\sigma_X^2=p(1-p)$\\
$\hat T^2=(N-2)\frac{R^2}{1-R^2}$ & $r_i=(Y_i-\Ybar)-\hat\beta_1(X_i-\Xbar)$ & $\SSR=\SST-\hat\beta_1^2S_{xx}$; divide by $\SST$\\
\bottomrule
\end{longtable}
}

\end{document}
